\chapter{Mathematics as the Language of Nature}
\label{chap:IV_language}

% =====================================================
% 4.1 Galileo Galilei and the New Physics
% =====================================================
Mathematics is not only a tool for describing nature,
but the form in which laws of nature can be formulated precisely in the first place.
Since the modern period, it has become increasingly clear
that physical knowledge arises where observation, measurement,
and mathematical structure come together.

This chapter follows this idea
from Galileo Galilei's new method
through Newton's universal laws
to functions, differential equations,
and modern physics.
In doing so, it becomes visible
that mathematics is not applied to nature afterward,
but makes its lawfulness recognizable in the first place.

\subsection{Galileo Galilei and the New Physics}

At the beginning of the modern period,
a radical change took place in the understanding of nature.
The world was no longer primarily \emph{interpreted},
but increasingly \emph{measured}.
This break is inseparably connected with the name
Galileo Galilei\index{Galilei, Galileo}.
He stands not only for new experiments or new instruments,
but for a new way of knowing nature at all.

Before Galileo Galilei, nature was usually described qualitatively.
Motion was regarded as something goal-directed,
heavy bodies were said to ``strive'' downward,
and heavenly bodies seemed to obey their own laws.
Explanations were verbal, often embedded in philosophy or theology.
Galileo Galilei opposed this with a different attitude:
natural phenomena should not primarily be interpreted,
but understood through their measurable relationships.

\begin{HistoryBox}[From Interpretation to Measurement]
	\phantomsection\label{box:measurement}
	The point of the ``new physics'' lies less in individual results
	than in the method:
	systematic observation, controlled measurement\index{measurement},
	and the conscious separation between \emph{description} and \emph{interpretation}.
	This shifted authority
	from tradition and argumentation
	to reproducible data.
\end{HistoryBox}

The decisive step, however, did not consist in measurement alone,
but in the mathematical formulation of what was measured.
Time\index{time}, distance\index{distance}, and velocity\index{velocity}
are not merely everyday words,
but abstract quantities
that enter into relationships with one another.
Galileo Galilei recognized:
only when motion is translated into numbers and relations
can general laws be recognized.

This fundamentally shifted the focus of physics.
The question was no longer:
\emph{Why does a body move in this way?}
Rather, it became:
\emph{Which mathematical relationship describes its motion?}

Here mathematics no longer serves merely as a calculation aid,
but becomes the medium of knowledge.
Laws of nature\index{law of nature}
no longer appear as verbal statements about hidden causes,
but as precise relationships between quantities.
Whether a body is made of wood or stone
is initially irrelevant for the mathematical law.
What matters is only the structure of the relationship.

Galileo Galilei's comparison of nature with a book
written in mathematical signs is therefore not a mere metaphor.
With it, he formulates a methodological insight:
whoever does not master the language of mathematics
cannot understand nature systematically.
This begins a development
that reaches from classical mechanics\index{mechanics}
to modern physics.

A particularly clear example of this change
is provided by Galileo Galilei's investigation of free fall\index{free fall}.
Since a falling body was too fast
to measure times precisely with the means available at the time,
Galileo Galilei artificially slowed down the motion:
he let balls roll down an inclined plane\index{inclined plane}.

This seemingly simple idea is decisive.
The motion remains physically similar,
but is stretched in time and thus becomes measurable.
Galileo Galilei determined the distance traveled
as a function of elapsed time
and discovered a surprising regularity:
the distance does not grow proportionally to time,
but proportionally to the square of time.
\[
s \propto t^{2}
\]

Thus motion was formulated for the first time
as a mathematical relationship\index{law of motion}.
The individual event is not decisive,
but the reproducible pattern.
Neither the material nor the mass of the ball is central,
but the structure of the dependence between space\index{space} and time.

This shows, in exemplary fashion,
what it means for mathematics to become the language of nature:
it forces observation into a form
that applies beyond the individual case.
Motion is no longer described only qualitatively,
but captured quantitatively.
Time appears as a physical quantity
that enters into a law.

Galileo Galilei thus provides not only individual results,
but a new principle of knowledge:
nature is understandable where it can be formulated mathematically.
After him, mathematics\index{mathematics}
is no longer merely a tool of physics,
but its foundation.

% =====================================================
% 4.2 Isaac Newton -- Mathematics Makes Nature Universal
% =====================================================
\subsection{Isaac Newton -- Mathematics Makes Nature Universal}

After Galileo Galilei, motion is mathematically describable.
At first, however, this description remains limited
to individual experiments and idealized situations.
The decisive next step consists in not only formulating
such relationships more precisely,
but extending them to all of nature.
This step is inseparably connected with the name
Isaac Newton\index{Newton, Isaac}.

Newton recognized
that the same mathematical relationships
that describe the fall of a body on Earth
must also apply to the motion of the heavenly bodies.
In doing so, he abolished a separation
that had lasted for millennia:
the distinction between earthly and heavenly physics.
For nature, there are no privileged places --
and for mathematics, none either.

\begin{HistoryBox}[Universal Laws Instead of Local Rules]
	\phantomsection\label{box:laws}
	What is new in Newton's physics is less the introduction of individual new phenomena
	than the claim to general validity.
	Laws of nature\index{law of nature} should not apply merely to certain situations,
	but without exception.
	Thus universality\index{universality}
	itself becomes a criterion of physical truth.
\end{HistoryBox}

Newton's true achievement therefore lies not only in individual formulas,
but in the insight
that laws of nature must be structural and universal.
Concepts such as force\index{force}, mass\index{mass},
and acceleration\index{acceleration}
are not intuitive properties,
but mathematically defined quantities.
Their meaning does not arise from everyday ideas,
but from the formal relationship
in which they occur.

This becomes especially clear in the law of gravitation\index{law of gravitation}.
A single mathematical relationship
describes both the fall of an apple
and the orbit of the planets\index{planetary motion}.
For the law, an apple is not a special case
and a planet is not a special object.
Both appear merely as mass points
within the same mathematical structure\index{structure}.

Thus a fundamental difference of earlier natural philosophy disappears.
The heavens do not follow their own laws,
and Earth is no exception.
What differs
are only the boundary conditions\index{boundary conditions},
not the form of the law.
The same mathematical structure captures both phenomena.

Newton's famous restraint toward metaphysical speculation --
\emph{hypotheses non fingo}\index{hypotheses non fingo}
(``I frame no hypotheses'') --
is, against this background, not a weakness,
but a methodological consequence.
Mathematics states
how nature behaves,
not why it exists in the ultimate metaphysical sense.
Precisely this limitation makes its statements universal.

With Newton, the mathematical description of nature
therefore reaches a new level.
Laws of nature are valid not merely
because they have often been confirmed,
but because their mathematical form
allows no exception.
The world is not calculable by accident;
it is calculable
because it is structured.

A particularly clear example of this approach
is Newton's second law of motion\index{law of motion}:
\[
F = m a
\]
This relationship is less an isolated measurement result
than a conceptual definition.
With it, Newton determines
what force is to mean in the first place:
force is that which produces acceleration.

This definition is deliberately abstract.
Force is not a visible cause,
not a material agent,
and not a directly experienced quantity.
It becomes accessible only through its effect
on the motion of a body.
Only through this mathematical relationship
does force become a precise physical concept.

The decisive step consists in introducing
acceleration as a fundamental quantity.
Not position\index{position},
not velocity\index{velocity},
but the change of velocity over time
becomes the measure of dynamical effect.
Thus the description of motion
detaches itself further from intuitive images
and becomes a structural relationship between quantities.

The law \(F = m a\) therefore does not explain
\emph{why} a body moves.
It determines
\emph{how} every interaction\index{interaction}
is to be described mathematically.
Precisely this formal restraint
makes the law universally applicable.
It holds regardless of whether the force is caused
by an impact,
by gravitation,
or by another interaction.

Newton's method thus shows, in exemplary fashion:
physical concepts receive their meaning
not through intuition,
but through their position within a mathematical relationship.
Mathematics\index{mathematics}
is here not merely the language of individual phenomena,
but the medium of universal description of nature.

% =====================================================
% 4.3 Functions and Laws of Motion
% =====================================================
\subsection{Functions and Laws of Motion}

With Newton, the mathematical description of nature reaches its universal form.
This, however, raises a further question:
In what mathematical form do laws of nature appear at all?
The answer does not lie in individual numerical values or tables,
but in functional relationships.
Laws of nature\index{law of nature} do not describe isolated states,
but dependencies between quantities.

A central example is the description of motion\index{motion}.
The position\index{position} of a body is not a fixed property,
but a quantity that changes with time\index{time}.
This description becomes physically meaningful only
when position is understood as a function\index{function} of time.
Motion is therefore not an object,
but a relationship between time and space\index{space}.

This functional point of view frees physics
from the intuition of individual moments.
What matters is not
where a body is located \emph{now},
but how its state changes with time.
Velocity\index{velocity}
and acceleration\index{acceleration}
are, in this sense, not additional things,
but derived concepts
that describe the time evolution of motion.

Laws of nature therefore do not simply appear
as finished trajectories.
Rather, they determine
how a quantity is allowed to change.
Mathematically, this leads to \emph{differential equations}\index{differential equation}.
A law of motion\index{law of motion}
does not directly connect position and time,
but relates rates of change to one another.
Only together with initial conditions\index{initial conditions}
does this yield a concrete motion.

Physics thus no longer describes only the \emph{result} of a motion,
but its \emph{rule}.
The future of a system is not given as a finished plan,
but arises from the local structure of the law.
Laws of nature therefore do not act in a goal-directed way,
but through instantaneous relationships.

This mathematical form reaches far beyond mechanics\index{mechanics}.
The same language is used
to describe growth\index{growth},
decay\index{decay},
or propagation\index{propagation}.
Whether a body falls,
a population grows,
or a charge is distributed in space --
in all cases, the issue is change
and its dependencies.
Different phenomena thus share the same basic mathematical form.

This is precisely why functions and differential equations
play such a central role in physics.
They are not chosen by chance,
but follow from the way
nature realizes processes.
Nature has no global plans,
but only local changes.
The differential equation represents exactly this structure.

From this perspective, classical mechanics appears
not as a special case,
but as a starting point.
Modern physics\index{modern physics}
adopts this mathematical language and generalizes it.
Fields\index{field},
wave functions\index{wave function},
and state spaces\index{state space}
replace particle trajectories,
yet the underlying principle remains:
laws of nature are functional relationships
that structure change.

A simple example makes this general form especially clear.
For uniformly accelerated motion\index{uniformly accelerated motion},
the velocity depends linearly on time:
\[
v(t) = a\,t,
\]
while the distance traveled grows quadratically with time:
\[
s(t) = \tfrac{1}{2} a\, t^{2}.
\]

These relationships are deliberately chosen to be simple.
They do not describe a single special situation,
but a general form of motion.
Whether the case is a falling body,
an accelerating vehicle,
or a technical system
is initially irrelevant for the mathematical form.
What matters is only
that the acceleration is constant.

This becomes especially intuitive in the graphical representation.
\begin{center}
	\begin{tikzpicture}[scale=1.1]
		
		% Axes
		\draw[->] (0,0) -- (4.5,0) node[right] {$t$};
		\draw[->] (0,0) -- (0,4.0) node[above] {$s(t)$};
		
		% Parabola s = 0.5 t^2 (form only, no numbers)
		\draw[thick, domain=0:4, samples=100]
		plot (\x, {0.25*\x*\x});
		
	\end{tikzpicture}
	
	\begin{tikzpicture}[scale=1.1]
		
		% Axes
		\draw[->] (0,0) -- (4.5,0) node[right] {$t$};
		\draw[->] (0,0) -- (0,4.0) node[above] {$v(t)$};
		
		% Straight line v = a t (form only, no numbers)
		\draw[thick] (0,0) -- (4,3);
		
	\end{tikzpicture}
\end{center}
\smallskip
\noindent
The upper curve shows the quadratic path \(s(t)\),
the lower straight line the linear behavior of the velocity \(v(t)\).

The velocity appears as a straight line in the time diagram,
the distance as a curved graph.
Both graphs are not additional assumptions,
but direct consequences of the same functional relationship.
They show
that laws of nature do not merely provide individual values,
but determine entire evolutions.

Precisely in this simple example,
it becomes visible what a law of motion accomplishes.
It does not describe one single path,
but a class of possible motions.
Different initial conditions change the evolution,
but not the underlying mathematical structure.

This demonstrates the generality of the functional description:
a single law generates infinitely many possible motions.
Physics does not describe the diversity of appearances,
but the structure\index{structure}
from which they arise.

% =====================================================
% 4.4 Differential Equations and Growth
% =====================================================

\subsection{Differential Equations and Growth}

To understand what a differential equation\index{differential equation} is,
a simple observation is enough at first:
nature is not static.
Everything we measure changes.
Positions\index{position} change with time\index{time},
temperatures rise or fall,
populations grow or shrink.
Physics\index{physics} therefore describes not only states,
but above all change\index{change}.

The term \emph{differential}\index{differential}
refers precisely to this idea of change.
It does not denote a new physical quantity,
but describes
how strongly a quantity changes over a very small interval.
A differential is therefore not an object,
but a mathematical form of local change.

A differential equation connects a quantity
with its own change.
It does not determine
which value a system has at a particular time,
but describes the rule
according to which it continues to develop
from its current state.
It therefore does not formulate a finished result,
but a prescription.

\begin{DidacticBox}[What a Differential Equation Describes]
	\phantomsection\label{box:differential}
	A differential equation does not prescribe a state
	or a finished solution.
	It determines
	how a system is allowed to change from its current state.
	
	It does not answer the question:
	\emph{``Where is the system?''}
	
	but:
	\emph{``How does it develop from here?''}
\end{DidacticBox}

A law of nature\index{law of nature} in this form
therefore does not act globally or in a goal-directed way.
It does not determine a final state,
but only the instantaneous development.
The future arises step by step
from the repeated application of the same local prescription.
Precisely this structure makes differential equations
so universally applicable.

This mathematical form is by no means
restricted to mechanics\index{mechanics}.
Its generality becomes especially intuitive
in the example of growth\index{growth}.
If one considers a system
whose rate of change is proportional to its current state,
then a simple
but far-reaching law results.
In this case, growth is
not the result of external planning,
but a direct consequence of local dynamics.

Mathematically, this means:
the change of a quantity depends on the quantity itself.
This structure leads to exponential growth
or exponential decay\index{decay}.
Whether the process involves biological populations,
radioactive processes,
or technical systems
is initially irrelevant for the mathematical form.
Different phenomena share the same equation.

Here in particular, the strength of mathematical description becomes visible.
The differential equation abstracts
from the concrete meaning of the quantities.
It knows neither animals nor atoms,
neither energy nor mass.
It describes only
a structural relationship between state and change.
The physical or biological interpretation
enters only in a second step.

This separation of mathematical form
and substantive meaning
is not a defect,
but a prerequisite for universality\index{universality}.
One and the same equation can describe growth,
decay, or stabilization\index{stabilization},
depending on parameters\index{parameter}
and initial conditions\index{initial conditions}.
The diversity of appearances therefore often does not arise
from entirely different laws,
but from different boundary conditions\index{boundary conditions}.

Differential equations thus make visible
what was already implicit in mechanics:
laws of nature are local prescriptions.
They do not operate globally,
but through immediate relationships.
Change is not a special case,
but the normal state of nature.

From this perspective, growth appears
not as something fundamentally different from motion,
but as a special form of dynamics\index{dynamics}.
Mathematics\index{mathematics}
does not distinguish between motion in space
and growth in time
when both follow the same structure.
Precisely there lies its power.

Differential equations and growth therefore mark
a further step
from intuitive mechanics
to abstract modern physics\index{modern physics}.
They show, in exemplary fashion,
how mathematics does not describe individual phenomena,
but determines the form of change itself.

% =====================================================
% 4.5 From Mechanics to Modern Physics
% =====================================================
\subsection{From Mechanics to Modern Physics}

Classical mechanics\index{mechanics} forms the historical starting point
of mathematical physics\index{physics}.
Its basic concepts are still intuitive:
particles\index{particle} have positions\index{position},
velocities\index{velocity}, and trajectories\index{trajectory};
forces\index{force} act between bodies;
and motion\index{motion} takes place in space\index{space}
and time\index{time}.
These ideas still shape our everyday physical understanding.

With the transition to modern physics\index{modern physics},
however, the mathematical basic principles do not change,
but rather the kind of objects being described.
The language remains,
but its range of application expands.
Particle trajectories are replaced by fields\index{field},
state functions\index{state function}, and abstract spaces.
Mechanics is not abolished,
but embedded in a larger framework.

This transition does not occur abruptly,
but consistently.
Already in mechanics, laws of motion\index{law of motion}
are formulated as differential equations\index{differential equation}.
Exactly the same basic mathematical structure later appears
in electrodynamics\index{electrodynamics},
in the theory of relativity\index{theory of relativity},
and in quantum mechanics\index{quantum mechanics}.
What changes
are the quantities being interpreted --
not the form of the laws.

In electrodynamics, effects no longer appear
as immediate forces between bodies,
but as fields
assigned to every point of space.
In relativity theory, space and time lose their absolute character
and themselves become part of the dynamical structure.
In quantum mechanics, finally,
the state of a system is no longer described by a trajectory,
but by a wave function\index{wave function}
that encodes probabilities\index{probability}.

\begin{DidacticBox}[Same Mathematics, New Concepts]
	\phantomsection\label{box:newconcepts}
	Modern physics does not replace the mathematical language of mechanics.
	It generalizes it.
	
	Differential equations, symmetries\index{symmetry},
	and conserved quantities\index{conserved quantity}
	remain.
	What changes
	is the meaning of the mathematical objects.
\end{DidacticBox}

The decisive point is:
increasing abstraction\index{abstraction}
does not mean a loss of reality,
but is a consequence of greater precision.
The more accurately natural phenomena are investigated,
the less intuitive images are able to carry the explanation.
Mathematics\index{mathematics}
takes over where imagination reaches its limits.

Modern physics thus shows, in a sharper form,
what was already implicit in mechanics:
laws of nature\index{law of nature}
are not stories about things,
but statements about structures\index{structure}.
They do not describe
what the world looks like,
but how change\index{change} is organized.
Mathematics is not a mere aid here,
but the actual medium of this description.

In this sense, the transition from mechanics to modern physics
does not mark a break,
but a consistent continuation.
The language remains the same,
but its scope grows.
What began with intuitive motion
becomes the abstract dynamics\index{dynamics}
of fields, states, and symmetries.

This closes the circle of this chapter.
From Galileo Galilei's measurements\index{Galilei, Galileo}
through Newton's universal laws\index{Newton, Isaac},
from functions\index{function}
and equations of motion
to differential equations and growth\index{growth},
one continuous principle becomes visible:
nature is understandable
where it is mathematically structured.
Modern physics is therefore not a countermodel to classical mechanics,
but its consistent next step.

\subsection*{Transition: From the Language of Nature to the Structure of Mathematics}

The examples so far show:
mathematics describes nature so successfully
not because it imitates reality,
but because it captures its relationships in abstracted form.
The more general a law of nature becomes,
the more strongly it detaches itself from the intuitive individual situation,
and the more clearly its structure emerges.

Precisely there lies the true strength of mathematics.
It does not cling to the surface of appearances,
but works with forms, relations, and patterns
that reach beyond the concrete case.
What appears in physics as a law
is, on the mathematical level, a structure.

Thus the language of nature leads directly
to a deeper insight:
mathematics gains its generality
precisely through abstraction.
Why abstraction does not mean distance from reality,
but rather access to its inner order,
will be shown in the next chapter.